\begin{lemma}[Mark data for a stopped contact walk]
Let $F$ be a finite-dimensional real inner-product space, let $C_0,\kappa,H\in\R$
and let $v\in\Aclass_{m,C_0}$ in the sense of \noderef{IsFitClass}. Let
$T\in\mathbb N$, $R\in\R$, $N\in\mathbb N$ and $\epsilon\in\R$, and let
$(\Omega,\mathcal M,\mu)$ be a measure space with a filtration
$(\mathcal G_n)_{n\in\mathbb N}$ carrying functions $\xi_n$, $z_n$, $z^+_n$,
$z^-_n$, $Q_0(n)$ and $\tau$ which form a stopped contact walk for
$(\kappa,H,T,R,v,N,\epsilon)$ in the sense of \noderef{IsContactWalk}. Assume
that for every $n$ the range of $z_n$ is at most countable. Then there are
functions $\sigma_j$, $p_j$, $a_j$ which are mark data for
$(C_0,v,N,\mathcal G,z,\tau)$ in the sense of \noderef{IsMarkData}.
\end{lemma}
\begin{proof}
Write $\delta_j=2^{-j}$ and $z_n=(x_n,t_n)$. Throughout, clause numbers refer to
\noderef{IsContactWalk}: clause 1 says $\mu(\Omega)=1$, clause 4 that each $z_n$
is $\mathcal G_n$-measurable, and clause 9 that $\tau$ is a stopping time for
$(\mathcal G_n)$.

\emph{Step 0: preliminaries.}
Since $\mu(\Omega)=1$ by clause 1, $\Omega$ is nonempty. Since each $z_n$ is
$\mathcal G_n$-measurable by clause 4 and the time coordinate $w\mapsto t(w)$ is
continuous, each function $\omega\mapsto t_n(\omega)$ is
$\mathcal G_n$-measurable, so
\begin{equation}\label{eq:cwm-level}
  \{t_m=\delta_j\}\in\mathcal G_m\qquad(m,j\in\mathbb N).
\end{equation}
Also, for every state $c$,
\begin{equation}\label{eq:cwm-eqset}
  \{\omega:z_i(\omega)=c\}\in\mathcal G_i ,
\end{equation}
because the set where two $\mathcal G_i$-measurable functions agree --- here
$z_i$ and the constant $c$ --- is $\mathcal G_i$-measurable. Finally, since
$\tau$ is a stopping time, $\{\tau\le n\}\in\mathcal G_n$ and
$\{\tau=k\}\in\mathcal G_k$ for every $k$, the latter because $\tau$ has
countable range; consequently, for $m\le n$,
\begin{equation}\label{eq:cwm-ge}
  \{m\le\tau\}=\Bigl(\bigcup_{k<m}\{\tau=k\}\Bigr)^{\rm c}\in\mathcal G_n ,
\end{equation}
a finite union of sets of $\mathcal G_k\subseteq\mathcal G_n$ with $k<m\le n$.

\emph{Step 1: the mark index and clauses 1--4 of \noderef{IsMarkData}.}
Fix $j$. If there is an $n\le\tau(\omega)$ with $t_n(\omega)=\delta_j$, let
$\sigma_j(\omega)$ be the least such $n$; otherwise put
$\sigma_j(\omega)=\tau(\omega)$. This is the index
$\tau\wedge\inf\{n:t_n=\delta_j\}$ of the paper, read as: a mark is the first
point on its time level, or the terminal point if the walk stops before that
level.

Then $\sigma_j\le\tau$: in the first case the least such $n$ is itself
$\le\tau(\omega)$, and in the second case $\sigma_j=\tau$. This is clause 2.

Let $n<\sigma_j(\omega)$. In the first case $n<\sigma_j(\omega)\le\tau(\omega)$,
so $n\le\tau(\omega)$, and minimality of $\sigma_j(\omega)$ rules out
$t_n(\omega)=\delta_j$; in the second case $n<\tau(\omega)$, and no $n'\le\tau$
at all has $t_{n'}=\delta_j$, in particular not $n$. Either way $n<\tau(\omega)$
and $t_n(\omega)\ne\delta_j$, which is clause 3.

In the first case $t_{\sigma_j(\omega)}(\omega)=\delta_j$ by the definition of
$\sigma_j(\omega)$, and in the second $\sigma_j(\omega)=\tau(\omega)$; this is
clause 4.

For clause 1 we claim
\begin{equation}\label{eq:cwm-stop}
  \{\sigma_j\le n\}
   =\{\tau\le n\}\cup\bigcup_{m\le n}\bigl(\{m\le\tau\}\cap\{t_m=\delta_j\}\bigr).
\end{equation}
For the inclusion $\subseteq$: if $\omega$ is in the first case, then
$m=\sigma_j(\omega)\le n$ satisfies $m\le\tau(\omega)$ and
$t_m(\omega)=\delta_j$, so $\omega$ lies in the union; if $\omega$ is in the
second case, then $\tau(\omega)=\sigma_j(\omega)\le n$. For $\supseteq$: if
$\tau(\omega)\le n$ then $\sigma_j(\omega)\le\tau(\omega)\le n$ by clause 2; and
if $m\le n$ satisfies $m\le\tau(\omega)$ and $t_m(\omega)=\delta_j$, then
$\omega$ is in the first case and $\sigma_j(\omega)\le m\le n$ by minimality.
Every set on the right of \eqref{eq:cwm-stop} lies in $\mathcal G_n$, by
$\{\tau\le n\}\in\mathcal G_n$, by \eqref{eq:cwm-ge} and by
\eqref{eq:cwm-level} together with $\mathcal G_m\subseteq\mathcal G_n$ for
$m\le n$. Hence $\sigma_j$ is a stopping time, which is clause 1. Since
$\sigma_j$ has countable range, it also follows that
\begin{equation}\label{eq:cwm-eqi}
  \{\sigma_j=i\}\in\mathcal G_i\qquad(i\in\mathbb N).
\end{equation}

\emph{Step 2: every function of the state at a fixed time is adapted.}
Fix $i$. By hypothesis the range of $z_i$ is countable, and it is nonempty
because $\Omega$ is, so there is a map $e:\mathbb N\to F\times\R$ whose range is
exactly the range of $z_i$. For $\omega\in\Omega$ let $\mathrm{cd}(\omega)$ be
the least $m$ with $e(m)=z_i(\omega)$; this exists because
$z_i(\omega)$ lies in the range of $e$. Then
$e(\mathrm{cd}(\omega))=z_i(\omega)$ for every $\omega$, and $\mathrm{cd}$ is
$\mathcal G_i$-measurable, because it takes values in the countable set
$\mathbb N$ and, for every $m$,
\[
  \{\mathrm{cd}=m\}
   =\{z_i=e(m)\}\cap\bigcap_{m'<m}\{z_i=e(m')\}^{\rm c}\in\mathcal G_i
\]
by \eqref{eq:cwm-eqset}. Now let $g$ be any function on states with values in a
finite-dimensional real inner-product space. For $k\in\mathbb N$ put
\[
  g_k(\omega)=\sum_{m<k}\mathbf 1_{\{\mathrm{cd}=m\}}(\omega)\,g(e(m)) ,
\]
a finite sum of indicators of sets of $\mathcal G_i$ times constants, hence
$\mathcal G_i$-measurable. For each $\omega$ and each $k>\mathrm{cd}(\omega)$
exactly one summand is nonzero, namely the one with $m=\mathrm{cd}(\omega)$, so
$g_k(\omega)=g(e(\mathrm{cd}(\omega)))=g(z_i(\omega))$; thus $g_k\to g\circ z_i$
pointwise, and a pointwise limit of $\mathcal G_i$-measurable functions with
values in a metrizable space is $\mathcal G_i$-measurable. Hence
\begin{equation}\label{eq:cwm-sel}
  \omega\mapsto g(z_i(\omega))\ \text{is }\mathcal G_i\text{-measurable}
  \qquad(i\in\mathbb N).
\end{equation}

\emph{Step 3: the fits and clauses 7--8.}
Fix $j$ and let $w$ be a state. Since $\delta_j>0$, the second clause of
\noderef{IsFitClass} provides a pair $(p(w),a(w))$ with
$\mathrm{IsFit}(C_0,v,w,\delta_j,p(w),a(w))$ in the sense of \noderef{IsFit};
fix one such pair for each $w$ (and each $j$). Define
\[
  G(w)=\begin{cases}(p(w),a(w)),&t(w)=\delta_j,\\(0,0),&t(w)\ne\delta_j,\end{cases}
  \qquad
  (p_j(\omega),a_j(\omega))=G\bigl(z_{\sigma_j(\omega)}(\omega)\bigr).
\]
If $t_{\sigma_j(\omega)}(\omega)=\delta_j$ then
$(p_j(\omega),a_j(\omega))=(p(w),a(w))$ at $w=z_{\sigma_j(\omega)}(\omega)$, so
$(p_j(\omega),a_j(\omega))$ is a fit for $v$ at that state of radius $\delta_j$,
which is clause 7; and if $t_{\sigma_j(\omega)}(\omega)\ne\delta_j$ then
$p_j(\omega)=0$ and $a_j(\omega)=0$, which is clause 8.

\emph{Step 4: clauses 5--6.}
Fix $j$ and $n$. We claim
\begin{equation}\label{eq:cwm-decomp}
  \mathbf 1_{\{\sigma_j\le n\}}\,G\bigl(z_{\sigma_j}\bigr)
   =\sum_{i\le n}\mathbf 1_{\{\sigma_j=i\}}\,G\bigl(z_i\bigr) .
\end{equation}
Indeed, at an $\omega$ with $\sigma_j(\omega)\le n$ exactly the summand
$i=\sigma_j(\omega)$ survives and equals
$G(z_{\sigma_j(\omega)}(\omega))$, while at an $\omega$ with
$\sigma_j(\omega)>n$ no summand survives and both sides are $0$. The sets
$\{\sigma_j=i\}$ with $i\le n$ are pairwise disjoint, lie in
$\mathcal G_i\subseteq\mathcal G_n$ by \eqref{eq:cwm-eqi}, and cover
$\{\sigma_j\le n\}$; these are the sets $A_{i,c}=\{\sigma_j=i\}\cap\{z_i=c\}$ of
the paper's countable tree of histories, grouped by $i$. By
\eqref{eq:cwm-sel} each $\omega\mapsto G(z_i(\omega))$ is
$\mathcal G_i$-measurable, hence $\mathcal G_n$-measurable for $i\le n$, so the
right-hand side of \eqref{eq:cwm-decomp} is a finite sum of products of
indicators of $\mathcal G_n$-measurable sets with $\mathcal G_n$-measurable
functions, and is therefore $\mathcal G_n$-measurable. Taking the two components
of \eqref{eq:cwm-decomp} --- the projections onto $F$ and onto $\R$ are
continuous --- gives that
$\omega\mapsto\mathbf 1_{\{\sigma_j\le n\}}(\omega)p_j(\omega)$ and
$\omega\mapsto\mathbf 1_{\{\sigma_j\le n\}}(\omega)a_j(\omega)$, that is the
functions equal to $p_j$ and $a_j$ on $\{\sigma_j\le n\}$ and to $0$ off it, are
$\mathcal G_n$-measurable. These are clauses 5 and 6.

All eight clauses of \noderef{IsMarkData} hold for every $j$ with
$1\le j\le N$ --- indeed for every $j$ --- so $(\sigma_j,p_j,a_j)$ is mark data
for $(C_0,v,N,\mathcal G,z,\tau)$. \end{proof}
